% संपूर्ण मराठी ग्रंथातील प्रकरण 50: कालिक तर्कशास्त्रे
% हा संपादनयोग्य स्रोतखंड आहे. संकलनासाठी संपूर्ण स्रोतसंग्रहातील
% अक्षररूपे, पूर्वभाग, चिन्हव्याख्या आणि आकृती-स्रोत आवश्यक आहेत.
% मूळ: Open Logic Project; मजकूर आणि रूपांतर CC BY 4.0.
% यंत्रानुवाद, दुरुस्ती आणि तपासणी: OpenAI Codex —
% GPT-5.6 Sol आणि GPT-6 Sol, दोन्ही Ultra effort.
% दुरुस्ती-पुनरावलोकन: GPT-6.1 Sol, Ultra effort.
\chapter{कालिक तर्कशास्त्रे}\label{aml:tl:chap}
\section{प्रस्तावना}\label{aml:tl:int:sec}
कालिक तर्कशास्त्रे यापुढे काय घडेल किंवा यापूर्वी
काय घडले आहे यांविषयीच्या विधानांचा विचार करतात.
आर्थर प्रायर यांना कालिक तर्कशास्त्राचे जनक मानले
जाते; त्यांनी त्याला \emph{काळवाचक तर्कशास्त्र}
असे नाव दिले. विधानीय तर्कशास्त्राची मूलभूत चौकट
विधाने साधारणतः काळवाचक भेद नसलेली मानते. त्या
चौकटीत कालिक संकारक घालण्याच्या प्रायर यांच्या
मूळ मोडल पद्धतीचे येथे आपण मुख्यतः अनुसरण करू.
उदाहरणार्थ, विधानीय तर्कशास्त्रात मी बीझी नावाच्या
कुत्र्याबद्दल बोलू शकतो. कुत्र्यांच्या स्वभावाप्रमाणे
ते कधी बसते आणि कधी बसत नाही. बीझी बसले आहे
आणि बीझी बसलेले नाही असे एकत्र म्हणणे अभिजात
तर्कशास्त्रात विरोधाभासी ठरेल. पण ही दोन्ही विधाने
खरी असू शकतात, फक्त एकाच वेळी नाहीत. भाषेत
कालिक संकारक घातल्यास ही गोष्ट तुलनेने सहज
व्यक्त करता येते. कालिक संकारकांमुळे ``बीझीला
खाऊ किंवा चेंडू मिळेल'' यावरून ``बीझीला खाऊ
मिळेल किंवा बीझीला चेंडू मिळेल'' यांसारख्या
अनुमानांची वैधताही स्पष्ट करता येते.
मात्र कालिक तर्कशास्त्रामुळे अनेक तत्त्वज्ञानविषयक
प्रश्न निर्माण होतात आणि त्यामुळे त्याची एक चौकट
निवडण्याऐवजी दुसरी निवडावी लागू शकते. उदाहरणार्थ,
भविष्यविषयक परिस्थितीवश विधान म्हणजे भविष्याबद्दलचे
असे विधान की जे अनिवार्यही नाही आणि अशक्यही नाही.
``रिचर्ड उद्या किराणा दुकानात जाईल'' असे म्हटल्यास,
आपण अजून न घडलेल्या घटनेविषयी विधान करतो; त्या
विधानाचे सत्यतामूल्य वादग्रस्त असते. किंबहुना,
त्याला मुळात सत्यतामूल्य \emph{देता} तरी येते का,
हाही वादाचा विषय आहे. आपण कठोर निर्धारणवादी
असू, तर संबंधित घटना घडण्यापूर्वीच हे विधान
खरे किंवा खोटे असते, एवढेच की त्याचे सत्यतामूल्य
आपल्याला अजून माहीत नसते, असे मानणे आपल्याला
मान्य होऊ शकते. याउलट, भविष्य खऱ्या अर्थाने
खुले आहे आणि भविष्यविषयक परिस्थितीवश विधानांची
सत्यतामूल्ये अजून अनिर्धारित आहेत, असेही आपण
मानू शकतो.
काळाची रचना आणि स्वरूप यांविषयीच्या अशा अनेक
भूमिका कालिक तर्कशास्त्रासाठी आपण निवडलेल्या
प्रतिमानांत आणि सैद्धांतिक चौकटींत अंतर्भूत असतात.
उदाहरणार्थ, काळ रेषीय, शाखायुक्त किंवा अगदी
वर्तुळाकार आहे अशी प्रतिमाने तयार करावीत का,
हा प्रश्न आपण विचारू शकतो. कालिक प्रतिमानांना
आरंभबिंदू आणि अंतबिंदू असावेत का, तसेच काळाचे
निरूपण विविक्त क्षणांनी करावे की सांतत्य म्हणून
करावे, हेही ठरवावे लागेल.
\section{कालिक तर्कशास्त्राचे अर्थनिर्धारण}\label{aml:tl:sem:sec}
\begin{defn}
कालिक तर्कशास्त्राच्या मूलभूत भाषेत पुढील गोष्टी असतात:
\begin{enumerate}
  \item असत्यता दर्शवणारा विधानीय स्थिरांक~$\lfalse$.
  \item सत्यता दर्शवणारा विधानीय स्थिरांक~$\ltrue$.
  \item विधानीय चलांचा गणनीय अनंत संच: $\Obj
    p_0$, $\Obj p_1$, $\Obj p_2$, \dots
  \item विधानीय संयोजके: \startycommalist
  \ycomma $\lnot$ (नकार)
  \ycomma $\land$ (संयोग)
  \ycomma $\lor$ (विकल्पयोग)
  \ycomma $\lif$ (शर्त)
  \ycomma $\liff$ (द्विशर्त).
  \item भूतकाळी संकारक $\Ptemp$ आणि $\Htemp$.
  \item भविष्यकाळी संकारक $\Ftemp$ आणि $\Gtemp$.
\end{enumerate}
\end{defn}
पुढे इतर प्रकारचे मोडल संकारक जोडण्याची शक्यताही पाहू.
\begin{defn}
कालिक भाषेतील \emph{सूत्रे} ची पुढीलप्रमाणे
  आगमनाने व्याख्या करतो:
\begin{enumerate}
\item $\lfalse$ हे अण्विक सूत्र आहे.
\item $\ltrue$ हे अण्विक सूत्र आहे.
\item प्रत्येक विधानीय चल~$\Obj p_i$ हे (अण्विक) सूत्र आहे.
\item $\varphi$ हे सूत्र असेल, तर $\lnot \varphi$
  हेही सूत्र आहे.
\item $\varphi$ आणि $\psi$ ही सूत्रे असतील, तर $(\varphi \land
  \psi)$ हे सूत्र आहे.
\item $\varphi$ आणि $\psi$ ही सूत्रे असतील, तर $(\varphi \lor \psi)$
  हे सूत्र आहे.
\item $\varphi$ आणि $\psi$ ही सूत्रे असतील, तर $(\varphi \lif \psi)$
  हे सूत्र आहे.
\item $\varphi$ आणि $\psi$ ही सूत्रे असतील, तर $(\varphi \liff \psi)$
  हे सूत्र आहे.
\item $\varphi$ हे सूत्र असेल, तर $\Ptemp  \varphi$, $\Htemp  \varphi$, $\Ftemp \varphi$, $\Gtemp  \varphi$
  ही सर्व सूत्रे आहेत.
\item याखेरीज दुसरे काहीही सूत्र नाही.
\end{enumerate}
\end{defn}
इतर आशयात्मक तर्कशास्त्रांप्रमाणे कालिक तर्कशास्त्रांचे
अर्थनिर्धारणही संबंधी प्रतिमानांच्या आधारे केले जाते.
\begin{defn}
  कालिक भाषेसाठीचे \emph{प्रतिमान} हे
  $\mModel{M} = \tuple{T, \prec, V}$ असे त्रिक असते. येथे
  \begin{enumerate}
  \item $T$ हा रिक्तेतर संच असून त्यातील घटकांचा अर्थ
    काळातील बिंदू असा लावला जातो.
  \item $\prec$ हा $T$ वरील द्विपदी संबंध आहे.
  \item $V$ हे प्रत्येक विधानीय चल~$p$ ला कालबिंदूंचा संच $V(p)$ देणारे फलन आहे.
  \end{enumerate}
  $t \prec t'$ असेल, तर $t$ हा~$t'$ च्या \emph{आधी येतो}
  असे म्हणतो. $t \in V(p)$ असेल, तर $p$ हे~$t$ येथे
  \emph{सत्य आहे} असे म्हणतो.
\end{defn}
सध्या आपण~$\prec$ या आधी येण्याच्या संबंधावर कोणतीही अट
घातलेली नाही, हे लक्षात येईल. त्यामुळे कालिक
प्रतिमानांच्या रचनेवर अजून कोणतेही निर्बंध नाहीत:
काळ रेषीय, शाखायुक्त, वर्तुळाकार किंवा इतर
कोणत्याही रचनेचा असलेली प्रतिमाने आपण घेऊ शकतो.
सामान्य मोडल तर्कशास्त्राप्रमाणे, प्रत्येक कालिक
प्रतिमानात कोणती सूत्रे कोणत्या बिंदूंवर
सत्य ठरतात हे निश्चित होते. संबंधी अर्थनिर्धारणाच्या
मूलभूत कल्पनेसाठी आपण तीच संज्ञा वापरतो:
``प्रतिमान~$\mModel{M}$ मध्ये सूत्र~$\varphi$
बिंदू~$t$ येथे सत्य ठरते.'' या संबंधाची आगमनाने
व्याख्या केली जाते; मोडल नसलेल्या सर्व संकारकांसाठी
ती सामान्य मोडल तर्कशास्त्रातील व्याख्येसारखीच आहे.
\begin{defn}\label{aml:tl:sem:defn:tmodels}
  प्रतिमान~$\mModel M$ मध्ये~$t$ येथे सूत्र~$\varphi$
  ची \emph{सत्यता}, म्हणजेच $\mSat{M}{\varphi}[t]$,
  पुढीलप्रमाणे आगमनाने परिभाषित होते:
  \begin{enumerate}
  \item \indcase{\varphi}{\lfalse}{$\mSat{M}{\lfalse}[t]$ कधीही नाही}.
  \item \indcase{\varphi}{\ltrue}{$\mSat{M}{\ltrue}[t]$ नेहमी असते}.
  \item $\mSat{M}{p}[t]$ जर आणि फक्त जर $t \in V(p)$
  \item \indcase{\varphi}{\lnot \psi}{$\mSat{M}{\indfrm}[t]$ जर आणि फक्त जर
    $\mSat/{M}{\psi}[t]$}.
  \item \indcase{\varphi}{(\psi \land \chi)}{$\mSat{M}{\indfrm}[t]$ जर आणि फक्त जर
    $\mSat{M}{\psi}[t]$ आणि $\mSat{M}{\chi}[t]$}.
  \item \indcase{\varphi}{(\psi \lor \chi)}{$\mSat{M}{\indfrm}[t]$ जर आणि फक्त जर
    $\mSat{M}{\psi}[t]$ किंवा $\mSat{M}{\chi}[t]$} (किंवा दोन्ही).
  \item \indcase{\varphi}{(\psi \lif \chi)}{$\mSat{M}{\indfrm}[t]$ जर आणि फक्त जर
    $\mSat/{M}{\psi}[t]$ किंवा $\mSat{M}{\chi}[t]$}.
  \item \indcase{\varphi}{(\psi \liff \chi)}{$\mSat{M}{\indfrm}[t]$ जर आणि फक्त जर
    एकतर $\mSat{M}{\psi}[t]$ आणि $\mSat{M}{\chi}[t]$ दोन्ही, किंवा
    $\mSat{M}{\psi}[t]$ आणि $\mSat{M}{\chi}[t]$ दोन्ही नाहीत}.
  \item\label{aml:tl:sem:defn:sub:mmodels-p}
    \indcase{\varphi}{\Ptemp  \psi}{$\mSat{M}{\indfrm}[t]$ जर आणि फक्त जर
    असा एखादा $t' \in T$ असेल की $t' \prec t$ आणि $\mSat{M}{\psi}[t']$}
\item\label{aml:tl:sem:defn:sub:mmodels-h}
    \indcase{\varphi}{\Htemp  \psi}{$\mSat{M}{\indfrm}[t]$ जर आणि फक्त जर
    $t' \prec t$ असलेल्या प्रत्येक $t' \in T$ साठी $\mSat{M}{\psi}[t']$}
   \item\label{aml:tl:sem:defn:sub:mmodels-f}
    \indcase{\varphi}{\Ftemp  \psi}{$\mSat{M}{\indfrm}[t]$ जर आणि फक्त जर
    असा एखादा $t' \in T$ असेल की $t \prec t'$ आणि $\mSat{M}{\psi}[t']$}
\item\label{aml:tl:sem:defn:sub:mmodels-g}
    \indcase{\varphi}{\Gtemp  \psi}{$\mSat{M}{\indfrm}[t]$ जर आणि फक्त जर
    $t \prec t'$ असलेल्या प्रत्येक $t' \in T$ साठी $\mSat{M}{\psi}[t']$}
  \end{enumerate}
\end{defn}
अर्थनिर्धारणावरून $\Ptemp$ आणि~$\Htemp$ हे
परस्पर द्वैती संकारक आहेत, तसेच $\Ftemp$ आणि
$\Gtemp$ हेही. म्हणून $\Htemp  \varphi$ ची व्याख्या
$\lnot \Ptemp \lnot \varphi$ अशी करता येईल; $\Gtemp$
आणि~$\Ftemp$ यांनाही हेच लागू होते.
\section{कालिक चौकटींचे गुणधर्म}\label{aml:tl:acc:sec}
आपली कालिक प्रतिमाने~$\prec$ या संबंधावर कोणतीही
अट घालत नसल्यामुळे, आपल्या परिचित स्वयंसिद्धकांपैकी
सर्व प्रतिमानांत लागू असणारे एकमेव स्वयंसिद्धक~$K$
आहे; किंवा त्याची $K_{\Gtemp}$ आणि~$K_{\Htemp}$
ही अनुरूप रूपे आहेत:
\begin{align*}
\tag{$K_{\Gtemp}$}  \Gtemp (p \to q) & \to (\Gtemp p \to \Gtemp q)\\
\tag{$K_{\Htemp}$}  \Htemp (p \to q) & \to (\Htemp p \to \Htemp q)
\end{align*}
मात्र प्रतिमानांनी काळाच्या निरूपणावर अधिक कडक
अटी घालाव्यात असे आपल्याला वाटत असेल—उदा.,
$\prec$ हा रेषीय क्रम असावा—तर त्या प्रतिमानांत
इतर सूत्रेही वैध ठरतील.
\begin{table}[t]
  \begin{tabular}{| p{.465\textwidth} || p{.465\textwidth} |}
    \hline
    {\emph{$\prec$ हा \dots असेल}} & {\emph{तर~$\mModel{M}$ मध्ये \dots सत्य आहे:}} \\
    \hline \hline
    \emph{संक्रमक}: \newline
    $\forall u \forall v \forall w ((u \prec v \land v \prec w) \lif u \prec w)$ &
    $\Ftemp \Ftemp p \lif \Ftemp p$  \\
    \hline
    \emph{रेषीय}: \newline
    $\forall w \forall v (w \prec v \lor w = v \lor v \prec w)$ &
    $(\Ftemp \Ptemp  p \lor \Ptemp \Ftemp p) \lif (\Ptemp p \lor p \lor \Ftemp p)$\\
    \hline
    \emph{सघन}: \newline
    $\forall w \forall v (w \prec v \to \exists u(w \prec u \land u \prec v))$ &
    $\Ftemp p \lif \Ftemp \Ftemp p$ \\
    \hline
    \emph{अपरिबद्ध (भूतकाळाकडे)}: \newline
    $\forall w \exists v( v \prec w)$ &
    $\Htemp p \to \Ptemp p$ \\
    \hline
    \emph{अपरिबद्ध (भविष्याकडे)}: \newline
    $\forall w \exists v( w \prec v)$ &
    $\Gtemp p \to \Ftemp p$ \\
    \hline
  \end{tabular}
  \caption{कालिक चौकटींचे काही अनुरूपता-गुणधर्म.}
  \label{aml:tl:acc:tab:correspondence}
\end{table}
\ref{aml:tl:acc:tab:correspondence} मधील अनेक गुणधर्म काळाचे
निरूपण करण्यासाठी बनवलेल्या प्रतिमानांत अपेक्षित
वाटू शकतात. तरीही, $\prec$ संबंधावर आपल्याला
हव्या त्या अटी घालता येत असल्या, तरी सर्वच अटी
कालिक तर्कशास्त्राच्या भाषेत व्यक्त करता येणाऱ्या
सूत्रे द्वारे दर्शवता येतात असे नाही. उदाहरणार्थ,
अपरावर्तीपणा, म्हणजे $\forall w \lnot (w \prec w)$
ही अट, कालिक तर्कशास्त्रातील कोणत्याही सूत्राशी
अनुरूप नाही.
\section{कालिक तर्कशास्त्रासाठी अतिरिक्त संकारक}\label{aml:tl:ext:sec}
भूतकाळ व भविष्यकाळासाठीच्या एकस्थानी संकारकांखेरीज,
कालिक तर्कशास्त्रात कधीकधी ``पासून'' आणि ``पर्यंत''
दर्शवणारे द्विस्थानी संकारक $\Since $ आणि~$\Until$
यांचाही समावेश असतो. म्हणजे कालिक तर्कशास्त्राच्या
भाषेत $\Since $ आणि~$\Until$ जोडून, कालिक
सूत्राच्या व्याख्येत पुढील नियम जोडायचा:
\begin{itemize}
\item[] $\varphi$ आणि $\psi$ ही सूत्रे असतील, तर
  $(\Since  \varphi \psi)$ आणि $(\Until \varphi \psi)$ ही दोन्ही
  सूत्रे आहेत.
\end{itemize}
या संकारकांचे अर्थनिर्धारण पुढीलप्रमाणे केले जाते:
\begin{defn}\label{aml:tl:ext:defn:since-until}
  प्रतिमान~$\mModel M$ मध्ये~$t$ येथे सूत्र~$\varphi$
  ची \emph{सत्यता}:
  \begin{enumerate}
  \item\label{aml:tl:ext:defn:sub:mmodels-since} \indcase{\varphi}{\Since  \psi
    \chi}{$\mSat{M}{\indfrm}[t]$ जर आणि फक्त जर असा एखादा $t' \in
    T$ असेल की $t' \prec t$ आणि $\mSat{M}{\psi}[t']$; तसेच
    $t' \prec s \prec t$ असलेल्या प्रत्येक $s$ साठी
    $\mSat{M}{\chi}[s]$}
   \item\label{aml:tl:ext:defn:sub:mmodels-until} \indcase{\varphi}{\Until \psi
    \chi}{$\mSat{M}{\indfrm}[t]$ जर आणि फक्त जर असा एखादा $t' \in
    T$ असेल की $t \prec t'$ आणि $\mSat{M}{\psi}[t']$; तसेच
    $t \prec s \prec t'$ असलेल्या प्रत्येक $s$ साठी
    $\mSat{M}{\chi}[s]$}
  \end{enumerate}
\end{defn}
$\Since  \psi \chi$ चा सहज अर्थ असा: ``ज्या
वेळेपासून~$\psi$ सत्य होते, त्या वेळेपासून~$\chi$
सत्य राहिले आहे.'' आणि $\Until \psi \chi$ चा सहज
अर्थ: ``ज्या वेळेपर्यंत~$\psi$ सत्य होईल, त्या
वेळेपर्यंत~$\chi$ सत्य राहील.''
\section{संभाव्य इतिहास}\label{aml:tl:poss:sec}
आतापर्यंत वापरलेली कालिक तर्कशास्त्राची संबंधी
प्रतिमाने अत्यंत लवचिक आहेत, कारण प्राप्यता संबंधावर
कोणतेही निर्बंध घालावे लागत नाहीत. त्यामुळे कालिक
प्रतिमाने भूतकाळात आणि भविष्यकाळात दोन्हीकडे
शाखायुक्त असू शकतात. पण शाखा फुटण्याची अधिक
``मोडल'' कल्पना आपण विचारात घेऊ शकतो: घटनांचे
अनुक्रम हे संभाव्य इतिहास मानायचे. यासाठी भाषा
बदलणे आवश्यकच नाही; तरीही आपण ``नेहमीचे''
मोडल संकारक $\Box$ आणि~$\Diamond$ जोडू शकतो.
कालांतराने कर्त्यांच्या ज्ञानात होणारे बदल दर्शवण्यासाठी
ज्ञानविषयक प्राप्यता संबंध जोडण्याचाही विचार करता येतो.
\begin{defn}
  कालिक भाषेसाठीचे \emph{संभाव्य इतिहासांचे प्रतिमान}
  हे $\mModel{M} = \tuple{T, C, V}$ असे त्रिक असते. येथे
  \begin{enumerate}
    \item $T$ हा रिक्तेतर संच असून त्यातील घटकांचा अर्थ
      काळातील अवस्था असा लावला जातो.
    \item $C$ हा प्रणालीच्या संगणकीय मार्गांचा, म्हणजे
      \emph{संभाव्य इतिहासांचा}, संच आहे. दुसऱ्या शब्दांत,
      $C$ हा $s_1$, $s_2$,~$s_3$, \dots अशा अवस्थांच्या
      $\sigma$ या अनुक्रमांचा संच आहे; येथे
      प्रत्येक $s_i \in T$. येथे $s_i$ या काळातील स्वतंत्र
      अवस्था-घटना मानतो: एका इतिहासात एकाच $T$-घटकाची
      पुनरावृत्ती होत नाही. प्रणालीची तीच स्थिती पुन्हा आली,
      तरी तिच्या वेगवेगळ्या काळातील घटनांना वेगळे घटक मानायचे.
      त्यामुळे $t,\sigma$ याने त्या इतिहासातील एकच स्थान
      निश्चित होते आणि $\prec_\sigma$ चा पुढील अर्थ संदिग्ध राहत नाही.
    \item $V$ हे प्रत्येक विधानीय चल~$p$ ला कालबिंदूंचा संच~$V(p) \subseteq T$ देणारे फलन आहे.
  \end{enumerate}
  गोष्टी सोप्या करण्यासाठी, एखादा इतिहास~$C$ मध्ये
  असेल, तर त्याचे सुरुवातीचे घटक वगळून मिळणारे
  सर्व उर्वरित अनुक्रमही त्यात आहेत, असे आपण साधारणतः
  मानू. उदाहरणार्थ, $s_1$, $s_2$,~$s_3$ हा अनुक्रम~$C$
  मध्ये असेल, तर $s_2$, $s_3$ आणि~$s_3$ हेही त्यात
  असतात. तसेच, अवस्था $s_i$ आणि~$s_j$ या
  अनुक्रम~$\sigma$ मध्ये असतील, तर $i < j$ असताना
  $s_i \prec_\sigma s_j$ असे म्हणतो.
  $t \in V(p)$ असेल, तर $p$ हे~$t$ येथे
  \emph{सत्य आहे} असे म्हणतो.
\end{defn}
येथे संबंधित बदल असा आहे: प्रतिमान~$\mModel M$
मध्ये कालबिंदू~$t$ येथे सूत्राची सत्यता
तपासताना आपण ती $t$ ही अवस्था-घटना असलेल्या
इतिहास~$\sigma \in C$ च्या सापेक्ष तपासतो.
विधानीय चलांची सत्यता फक्त $t$ वर अवलंबून असते:
$\mSat{M}{p}[t,\sigma]$ जर आणि फक्त जर $t\in V(p)$.
सत्यता-फलनात्मक संयोजकांच्या अटी
\ref{aml:tl:sem:defn:tmodels} मधीलप्रमाणेच राहतात;
मात्र त्यांतील प्रत्येक उपसूत्राची सत्यता त्याच
$t,\sigma$ येथे तपासायची. उदाहरणार्थ,
$\mSat{M}{\lnot \psi}[t,\sigma]$ जर आणि फक्त जर
$\mSat/{M}{\psi}[t,\sigma]$; आणि
$\mSat{M}{\psi\land \chi}[t,\sigma]$ जर आणि फक्त जर
$\mSat{M}{\psi}[t,\sigma]$ आणि $\mSat{M}{\chi}[t,\sigma]$.
संयोजक स्वतः इतिहास बदलत नसले, तरी त्यांची कालिक किंवा
मोडल उपसूत्रे त्या इतिहासावर अवलंबून असू शकतात.
आता या इतिहासांच्या सापेक्ष आपण
भविष्यकाळी संकारक~$\Ftemp$ ची पुनर्व्याख्या करतो
आणि~$\Diamond$ संकारक जोडतो.
\begin{defn}\label{aml:tl:poss:defn:phmodels}
  प्रतिमान~$\mModel M$ मध्ये~$t, \sigma$ येथे
  सूत्र~$\varphi$ ची \emph{सत्यता}, म्हणजेच
  $\mSat{M}{\varphi}[t, \sigma]$:
  \begin{enumerate}
  \item\label{aml:tl:poss:defn:sub:phmodels-f} \indcase{\varphi}{\Ftemp
    \psi}{$\mSat{M}{\indfrm}[t, \sigma]$ जर आणि फक्त जर
    असा एखादा $t' \in T$ असेल की $t \prec_\sigma t'$
    आणि $\mSat{M}{\psi}[t', \sigma]$.}
  \item\label{aml:tl:poss:defn:sub:phmodels-diamond} \indcase{\varphi}{\Diamond
    \psi}{$\mSat{M}{\indfrm}[t, \sigma]$ जर आणि फक्त जर
    असा एखादा $\sigma' \in C$ असेल की त्यात $t$ येतो
    आणि $\mSat{M}{\psi}[t, \sigma']$.}
  \end{enumerate}
\end{defn}
इतर कालिक आणि मोडल संकारकांचीही अशीच व्याख्या
करता येते. आता आपण काळवाचकता आणि मोडलता
एकत्र करणारी विधाने दर्शवू शकतो. उदाहरणार्थ,
``$p$ पुढे घडणार नाही, पण घडू शकला असता'' हे
विधान $\lnot \Ftemp p \land \Diamond \Ftemp p$
या सूत्राने दर्शवता येईल. ज्या बिंदू $t$ आणि इतिहास $\sigma$ येथे
त्या इतिहासातील कोणत्याही पुढील अवस्थेत $p$ सत्य नसते,
पण $t$ असलेल्या एखाद्या पर्यायी इतिहासात $p$ पुढील
एखाद्या अवस्थेत सत्य होते, तेथे हे सूत्र सत्य ठरेल.
